\section{Discussão}\label{chap:discussao-final}

\subsection{Atendimento aos objetivos e evolução do escopo}

Os objetivos de integração, seleção de áreas, índices, mapas, séries e exportação foram materializados em módulos verificáveis. A \cref{tab:objetivos-final} distingue implementação e extensão da avaliação, sem equiparar testes de software a estudos agronômicos.

\begin{table*}[htbp]
\centering\small
\caption{Síntese do atendimento aos objetivos da redação final.}
\label{tab:objetivos-final}
\begin{tabular}{p{0.10\textwidth}p{0.20\textwidth}p{0.61\textwidth}}
\hline
\textbf{Objetivo} & \textbf{Situação} & \textbf{Fundamento e ressalva} \\ \hline
OE1 & Implementado & Estrutura modular e métodos documentados; bibliografia com pendências explícitas. \\
OE2 & Implementado & Busca, desenho, gravação e recuperação; sessão web não é persistência permanente. \\
OE3 & Implementado & Filtros e fórmulas exercitados, demonstração remota executada; adequação histórica da máscara requer avaliação específica. \\
OE4 & Implementado & Composição, média e consulta local; contrato de escala corrigido e testado. \\
OE5 & Implementado & Série, tendência e regra unilateral exercitadas com dados controlados; sem acurácia agronômica atribuída. \\
OE6 & Implementado & Clima compartilhado, unidade normalizada e resumo propagado; ausência tratada explicitamente. \\
OE7 & Implementado & PDF, XLSX, CSV ZIP e GeoJSON exercitados com entradas controladas; não equivalem a pacote completo de reprodução. \\
OE8 & Parcialmente concluído & Suíte e fluxo real com downloads executados; pesquisa com usuários sem resultados e estudos por cultura não realizados. \\ \hline
\end{tabular}
\end{table*}

Na proposta, NDVI e NDWI eram o foco e NDMI aparecia como evolução futura. A implementação incorporou NDMI, exploração numérica no navegador, modos de armazenamento e exportações estruturadas. A ausência original de banco foi substituída por SQLite local e, posteriormente, por alternativa de sessão para o ponto de entrada web. Essas adições são resultados de engenharia e não substituem o objetivo original de validar o protótipo em áreas de soja, milho e pastagem.

Os estudos de soja, milho e pastagem, experimentos controlados e coleta em campo previstos não foram realizados. A avaliação concentrou-se no artefato computacional. Essa delimitação mantém a contribuição de engenharia, mas não substitui o objetivo experimental original e requer apreciação do orientador.

Outros requisitos das especificações também exigem distinção. O fluxo atual mostra área em hectares, sem seletor de acres, e não implementa bloqueio/liberação explícitos de navegação como descrito em especificações anteriores. A exclusão de localizações está no serviço, mas não na interface. Descrever esses itens como integralmente disponíveis apenas porque constam no planejamento produziria uma representação inadequada do sistema.

\subsection{Contribuições de Engenharia de Software}

A plataforma articula representações heterogêneas: texto de busca, coordenadas, polígonos, operações remotas de imagem, matrizes comprimidas, séries e arquivos. O encapsulamento em serviço e repositórios permite manter as regras de nome e transferência independentemente do armazenamento local ou de sessão. O snapshot de exportação reduz a possibilidade de gerar arquivos associados a parâmetros diferentes dos que produziram os dados.

O componente de consulta local é uma decisão de interação relevante: transfere uma grade após a análise para evitar requisições por movimento do cursor. Em contrapartida, impõe custo de memória e transferência, depende de suporte do navegador e limita o tamanho da região. Trata-se de um compromisso explícito de engenharia, não de comprovação de baixa latência em qualquer área.

As exportações disponibilizam informações fora da interface e permitem conferir dados com outras ferramentas. O PDF sintetiza, enquanto tabelas e GeoJSON preservam estruturas reutilizáveis. Contudo, não há exportação de toda a proveniência necessária à reconstrução exata do processamento. Interoperabilidade e transparência são ampliadas, mas reprodutibilidade integral ainda requer versões, IDs, parâmetros e dados auxiliares adicionais.

\subsection{Consistência das integrações e limites dos testes}

A divergência corrigida na radiação evidencia a diferença entre cobertura e correção semântica: a execução dos caminhos não verifica, por si só, a unidade de origem. A interpretação explícita dos metadados e sua reutilização em resumo e exportações reduziram a duplicação dessa regra.

A legenda passou a compartilhar o intervalo dos tiles. O mapa interativo favorece contraste interno ao recorte, enquanto o PDF mantém escala fixa, útil para referência comum entre relatórios. Aparências distintas podem ser corretas quando cada escala é comunicada.

A harmonização dos pipelines e a associação por aquisição preservam a interpretação dos registros entre modos e formatos. O identificador melhora a proveniência tabular, embora ainda não constitua inventário completo de processamento. A invalidação de estado complementa esse cuidado ao impedir apresentação de resultados sob parâmetros analíticos diferentes.

\subsection{Limitações espaciais, temporais e estatísticas}

A máscara QA60 não resolve todos os problemas de nuvens e sombras e possui restrições históricas de disponibilidade. O filtro de nuvens é do produto, não da área. A grade a 10 m não confere resolução nativa de 10 m às bandas B8A e B11. A média espacial do pipeline não fixa explicitamente o CRS, enquanto a consulta interativa utiliza UTM; correspondência numérica estrita requer verificação adicional.

A composição mediana pode reduzir a influência de observações extremas, mas também oculta quando uma mudança ocorreu e quantas observações sustentam cada posição. A série pode apresentar lacunas, datas repetidas e variação do suporte espacial válido. O resumo de imagens não informa essa cobertura por pixel.

A heurística de anomalias considera apenas três observações e desvios negativos. Não sinaliza quedas após janela de desvio zero, não modela sazonalidade ou fenologia e não separa ruído de mudanças reais. A tendência usa ordem de observações, não tempo decorrido. Sem referências agronômicas e avaliação apropriada, esses produtos servem à exploração de comportamentos atípicos e não ao diagnóstico automático.

\subsection{Operação, serviços e avaliação não realizada}

O uso de bibliotecas abertas deve ser distinguido do acesso a serviços externos. Earth Engine e Esri não são componentes locais de código aberto da aplicação. Dados Sentinel possuem condições próprias de uso, o serviço de geocodificação público possui política de utilização e o Earth Engine requer projeto e elegibilidade, com cotas e condições específicas \citep{GOOGLE_EE_NONCOMMERCIAL,NOMINATIM_POLICY}. Não se demonstrou operação comercial gratuita ou custo total inferior ao de outras soluções.

O serviço público Nominatim limita o tráfego agregado da aplicação a uma requisição por segundo e proíbe autocomplete. O código utiliza busca explícita e identifica a aplicação, mas não implementa limitador global ou cache de buscas. A atribuição e as condições de acesso devem ser revistas antes de ampliar o número de visitantes; disponibilidade pública não implica capacidade ilimitada.

Na web, visitantes compartilham a conta de serviço e o consumo de cotas, embora os registros estejam separados por sessão. GeoJSON e arquivos de resultados podem conter coordenadas sensíveis, e a privacidade dos arquivos transferidos fica sob responsabilidade de quem os armazena e compartilha. O isolamento de sessão não constitui sistema de autorização de propriedades.

Não foram conduzidos testes de carga, ensaios representativos de latência, comparação quantitativa com plataformas ou avaliação de campo. A pesquisa de experiência de uso permanece sem resultados analisados; sua amostragem por adesão limitará generalizações. A resolução das bandas e dos dados meteorológicos, a dependência de serviços e cotas, a memória da grade no navegador e a natureza heurística das anomalias delimitam o uso. A máscara QA60 foi preservada: uma substituição por SCL ou probabilidade de nuvens exige verificação de classes, sombras, disponibilidade histórica e efeitos sobre os resultados, não apenas troca de API.

A inspeção visual identificou contraste insuficiente em alguns textos quando os estilos escuros se combinam ao tema de widgets do ambiente. A emulação móvel confirmou reorganização e lookup por toque, mas não substitui ensaios de acessibilidade ou uma matriz de dispositivos e navegadores.
